\chapter{\ifenglish Introduction\else บทนำ\fi}

\section{\ifenglish Project rationale\else ที่มาของโครงงาน\fi}

ระบบ Security Information and Event Management (SIEM) ทำหน้าที่รวบรวม Security Logs
จากแหล่งข้อมูลจำนวนมากในองค์กร เช่น เครื่องลูกข่าย เซิร์ฟเวอร์ และอุปกรณ์เครือข่าย
แล้วนำมาจัดเก็บ ค้นหา และสร้างการแจ้งเตือน (Alert) เมื่อพบพฤติกรรมที่น่าสงสัย
อย่างไรก็ดี ปริมาณและความหลากหลายของ Logs ทำให้การสืบสวนเหตุการณ์ทำได้ยาก
ตัวอย่างเช่น ชุดข้อมูลจำลองการโจมตีแบบ Account Discovery (T1087.001)
จาก Splunk Attack Data~\cite{splunk-attack-data} มี Windows Sysmon Logs ถึง 11,090 events
ที่เกิดขึ้นภายในเวลาประมาณ 15 นาที (ขนาดไฟล์ XML ประมาณ 16.8 MB)
ซึ่งมีเพียงส่วนน้อยเท่านั้นที่เกี่ยวข้องกับการโจมตีจริง

แม้ Alert จาก SIEM หรือ Endpoint Detection and Response (EDR) จะช่วยชี้ Event หรือกิจกรรมที่ควรตรวจสอบ
แต่การระบุว่าเหตุการณ์ดังกล่าวเป็นการโจมตีหรือไม่ และสอดคล้องกับเทคนิคใดใน MITRE ATT\&CK~\cite{mitre-attack}
ยังต้องอาศัย SOC Analyst ในการรวบรวมและเชื่อมโยง Logs ที่เกี่ยวข้องเข้าด้วยกัน
ซึ่งใช้เวลาและความเชี่ยวชาญสูง

แนวทางหนึ่งที่ได้รับความสนใจคือการใช้โมเดลภาษาขนาดใหญ่ (Large Language Model: LLM) ช่วยวิเคราะห์ Logs
แต่การส่ง Raw Logs ให้ LLM โดยตรงมีปัญหาสำคัญ 2 ประการ ได้แก่
(1) ต้นทุนด้าน Token สูง เนื่องจาก Logs มีขนาดใหญ่และมีโครงสร้างซ้ำซ้อน และ
(2) Logs มี Noise จำนวนมาก (เหตุการณ์ที่ไม่เกี่ยวข้องกับการวิเคราะห์ รวมถึง Field และ Event ที่ไม่จำเป็น)
ซึ่งทำให้ LLM ตีความผิดพลาดได้ง่าย
จากการทดลองเบื้องต้นของคณะผู้จัดทำ เมื่อส่ง Raw Logs ของชุดข้อมูล T1087.001 ที่ดึงจาก Splunk ผ่าน API
ให้ LLM เชิงพาณิชย์หลายตัววิเคราะห์โดยตรง พบว่าบางโมเดลระบุเทคนิคผิด
(เช่น ตอบเป็น T1003 OS Credential Dumping) บางโมเดลสรุปว่าไม่พบการโจมตี
และแม้โมเดลที่ตอบใกล้เคียงที่สุดก็ระบุได้เพียงระดับ Technique (T1087) แต่ไม่ถึงระดับ Sub-technique ที่ถูกต้อง

ด้วยเหตุนี้ โครงงานนี้จึงมุ่งออกแบบและประเมิน \textbf{Context Construction Pipeline}
ที่ทำหน้าที่สร้าง \textbf{Investigation Context} ซึ่งหมายถึงชุดเหตุการณ์และความรู้ที่เกี่ยวข้องกับเหตุการณ์ที่น่าสงสัย
ที่ถูกจัดเตรียมไว้เพื่อให้ LLM ใช้เหตุผลกับเหตุการณ์นั้น
โดย Pipeline จะนำ Alert เป็นจุดเริ่มต้นการสืบสวน แล้วทำการ Normalize, คัดเลือก Event ที่เกี่ยวข้อง (Event Selection)
และแปลงรูปแบบข้อมูล (Representation) ให้กระชับ ก่อนเสริมด้วยองค์ความรู้จาก MITRE ATT\&CK
ผ่าน Retrieval-Augmented Generation (RAG)~\cite{lewis2020rag}
เพื่อให้ LLM ขนาดเล็กที่ทำงานภายในเครื่อง (Local LLM) สามารถระบุเทคนิคการโจมตีได้อย่างแม่นยำ
โดยไม่ต้องส่งข้อมูล Logs ขององค์กรออกไปยังบริการภายนอก

\section{\ifenglish Objectives\else วัตถุประสงค์ของโครงงาน\fi}
\begin{enumerate}
    \item เพื่อออกแบบและพัฒนา Context Construction Pipeline สำหรับสร้าง Investigation Context
          จาก Security Logs ในระบบ SIEM ก่อนส่งให้ LLM วิเคราะห์
    \item เพื่อเปรียบเทียบประสิทธิภาพของวิธีการคัดเลือก Event (Event Selection)
          และรูปแบบการนำเสนอข้อมูล (Representation) แบบต่าง ๆ
    \item เพื่อทดสอบว่า Investigation Context ช่วยเพิ่มความแม่นยำในการระบุเทคนิค MITRE ATT\&CK
          ได้ดีกว่าการใช้ Raw Logs
    \item เพื่อลดปริมาณข้อมูลที่ส่งให้ LLM (ลด Token Cost และลด Noise)
          โดยยังคงหลักฐานที่จำเป็นต่อการวิเคราะห์ไว้ครบถ้วน
    \item เพื่อพัฒนาต้นแบบ Web Dashboard ที่นำเสนอผลการวิเคราะห์เป็นรายงานการสืบสวน (Investigation Report)
          ให้แก่ SOC Analyst
\end{enumerate}

\section{\ifenglish Project scope\else ขอบเขตของโครงงาน\fi}

\subsection{\ifenglish Hardware scope\else ขอบเขตด้านฮาร์ดแวร์\fi}
\begin{enumerate}
    \item ระบบทำงานบนเครื่องคอมพิวเตอร์ส่วนบุคคลหรือเซิร์ฟเวอร์ที่มีหน่วยประมวลผลกราฟิก (GPU)
          เพียงพอสำหรับรัน LLM ขนาดประมาณ 9 พันล้านพารามิเตอร์ (9B) ภายในเครื่อง
    \item ไม่มีการพัฒนาหรือติดตั้งอุปกรณ์ฮาร์ดแวร์เฉพาะทาง เช่น เซนเซอร์หรืออุปกรณ์เครือข่าย
\end{enumerate}

\subsection{\ifenglish Software scope\else ขอบเขตด้านซอฟต์แวร์\fi}
\begin{enumerate}
    \item ใช้ชุดข้อมูลสาธารณะที่มี Ground Truth เพื่อให้ตรวจสอบความแม่นยำได้
          ได้แก่ OTRF Security-Datasets~\cite{otrf-security-datasets} สำหรับออกแบบและวัดผล (Benchmark)
          และ Splunk Attack Data~\cite{splunk-attack-data} สำหรับทดสอบ (Validation)
          โดยนำ Logs ไป Replay เข้า Splunk Enterprise เพื่อจำลองการทำงานของ SIEM จริง
    \item รองรับ Logs ในรูปแบบ Windows Sysmon~\cite{sysmon} และ Linux Auditd~\cite{auditd}
    \item หน่วยของการวิเคราะห์คือ 1 Alert ต่อ 1 กรณีการสืบสวน (Alert-centric)
          โดย Alert ถือเป็น Input ที่ได้รับจาก SIEM/EDR เดิม โครงงานนี้ไม่ได้ทำหน้าที่ตรวจจับภัยคุกคามแทน SIEM/EDR
    \item ใช้ LLM ตัวเดียวตลอดการทดลอง (Qwen3.5 9B ทำงานภายในเครื่องผ่าน Ollama)
          เพื่อให้ Pipeline เป็นตัวแปรเดียวที่เปลี่ยนแปลงในการทดลอง
    \item Ground Truth อยู่ในระดับ Scenario Dataset กล่าวคือ 1 Scenario Dataset มีเทคนิค ATT\&CK
          ที่เป็นเฉลยหนึ่งเทคนิค ซึ่งใช้สำหรับประเมินผลเท่านั้น ไม่ได้นำไปใช้ฝึกสอนโมเดล
    \item ผลลัพธ์ของระบบคือเทคนิค MITRE ATT\&CK ที่ระบุได้ พร้อมหลักฐานและเหตุผล
          นำเสนอเป็น Investigation Report ผ่าน Web Dashboard
\end{enumerate}

\section{\ifenglish Expected outcomes\else ประโยชน์ที่ได้รับ\fi}
\begin{enumerate}
    \item ได้ Context Construction Pipeline ที่สามารถลดปริมาณข้อมูลและ Token ที่ส่งให้ LLM
          โดยยังคงหลักฐานที่จำเป็นต่อการระบุเทคนิคการโจมตี
    \item ได้ผลการเปรียบเทียบเชิงทดลองของวิธี Event Selection และ Representation แบบต่าง ๆ
          รวมถึงผลของการใช้ RAG ต่อความแม่นยำของ LLM ขนาดเล็ก
    \item ได้ระบบต้นแบบ LLM-based SIEM Investigation Assistant ที่เชื่อมต่อกับ SIEM (Splunk)
          และสร้าง Investigation Report ให้ SOC Analyst ผ่าน Web Dashboard
    \item ช่วยให้ SOC Analyst ลดเวลาในการสืบสวนเหตุการณ์ และช่วยให้องค์กรลดต้นทุนการใช้ LLM
          โดยสามารถใช้โมเดลที่ทำงานภายในองค์กรได้
    \item คณะผู้จัดทำได้เรียนรู้และประยุกต์ใช้ความรู้ด้านความมั่นคงปลอดภัยไซเบอร์ ระบบ SIEM
          และการประยุกต์ใช้ LLM กับข้อมูลจริง
\end{enumerate}

\section{\ifenglish Technology and tools\else เทคโนโลยีและเครื่องมือที่ใช้\fi}

\subsection{\ifenglish Hardware technology\else เทคโนโลยีด้านฮาร์ดแวร์\fi}
\begin{enumerate}
    \item เครื่องคอมพิวเตอร์ส่วนบุคคลที่มี GPU สำหรับรัน LLM และ Embedding Model ภายในเครื่อง
          % TODO: ระบุสเปกเครื่อง เช่น CPU, RAM, รุ่น GPU และขนาด VRAM
    \item เครื่องเสมือนหรือเซิร์ฟเวอร์สำหรับติดตั้ง Splunk Enterprise เพื่อจำลองระบบ SIEM
\end{enumerate}

\subsection{\ifenglish Software technology\else เทคโนโลยีด้านซอฟต์แวร์\fi}
\begin{enumerate}
    \item \textbf{Splunk Enterprise} -- ระบบ SIEM สำหรับรับและจัดเก็บ Logs ที่ Replay จากชุดข้อมูล
          และเป็นแหล่งข้อมูลที่ Pipeline ดึง Logs ผ่าน REST API
    \item \textbf{Python} -- ภาษาหลักสำหรับพัฒนา Pipeline ได้แก่ การ Parse/Normalize Logs,
          Event Selection, Representation, RAG และการประเมินผล
    \item \textbf{Ollama}~\cite{ollama} -- เครื่องมือสำหรับรัน LLM ภายในเครื่อง
    \item \textbf{Qwen3.5:9B} -- LLM ที่ใช้วิเคราะห์ Investigation Context
          และ \textbf{qwen3-embedding:0.6b} สำหรับสร้าง Embedding ของฐานความรู้ใน RAG
    \item \textbf{Hermes} -- Harness สำหรับจัดการ Prompt และการเรียกใช้งานโมเดล
    \item \textbf{MITRE ATT\&CK Enterprise v19.2 (STIX)} -- ฐานความรู้เทคนิคการโจมตีที่ใช้สร้าง Knowledge Base ของ RAG
    \item \textbf{React (TypeScript) + Vite} -- พัฒนาส่วนหน้า (Frontend) ของ Web Dashboard
    \item \textbf{Git/GitHub} และ \textbf{Visual Studio Code} -- จัดการซอร์สโค้ดและพัฒนาโปรแกรม
\end{enumerate}

\section{\ifenglish Project plan\else แผนการดำเนินงาน\fi}

แผนการดำเนินงานของโครงงานแสดงดังตารางด้านล่าง
ณ วันที่จัดทำรายงานนี้ ขั้นตอนการตกลงหัวข้อ การรวบรวมข้อมูล การออกแบบภาพรวม Pipeline
และการติดตั้ง Splunk พร้อมป้อนข้อมูลจำลองได้ดำเนินการเสร็จสิ้นแล้ว
ส่วนการออกแบบ Event Selection \& Representation การทำ Ground Truth
การพัฒนา RAG และ Web Dashboard อยู่ระหว่างดำเนินการ

\begin{plan}{6}{2026}{3}{2027}
    \planitem{6}{2026}{6}{2026}{ตกลงหัวข้อโครงงาน}
    \planitem{7}{2026}{8}{2026}{รวบรวมข้อมูล (ชุดข้อมูลและงานที่เกี่ยวข้อง)}
    \planitem{7}{2026}{7}{2026}{ออกแบบภาพรวม Pipeline}
    \planitem{8}{2026}{8}{2026}{ติดตั้ง Splunk และป้อนข้อมูลจำลอง}
    \planitem{9}{2026}{9}{2026}{ออกแบบ Event Selection \& Representation}
    \planitem{10}{2026}{10}{2026}{จัดทำ Ground Truth}
    \planitem{10}{2026}{10}{2026}{สร้าง Parser และ Backend API}
    \planitem{11}{2026}{11}{2026}{Implement Event Selection \& Representation}
    \planitem{11}{2026}{12}{2026}{พัฒนา RAG}
    \planitem{12}{2026}{12}{2026}{เขียน System Prompt}
    \planitem{12}{2026}{1}{2027}{พัฒนา Web Dashboard}
    \planitem{1}{2027}{1}{2027}{Integration Testing}
    \planitem{1}{2027}{2}{2027}{การทดลองและการประเมินผล}
    \planitem{2}{2027}{3}{2027}{จัดทำรูปเล่มรายงาน}
\end{plan}

\section{\ifenglish Roles and responsibilities\else บทบาทและความรับผิดชอบ\fi}

คณะผู้จัดทำแบ่งบทบาทและหน้าที่ตามความถนัด ดังนี้
\begin{enumerate}
    \item \textbf{นายธนกร เสาคำ} (Full-Stack \& Validation) --
          พัฒนา Web Dashboard ทั้งส่วน Frontend และ Backend API
          เชื่อมต่อระบบกับ Splunk และตรวจสอบความถูกต้องของผลลัพธ์ (Validation)
          ใช้ความรู้ด้าน Software Engineering, Web Development และฐานข้อมูล
    \item \textbf{นางสาวนพรดา กฤษวัฒนานนท์} (AI \& Experiment) --
          พัฒนาส่วน RAG และ LLM Analysis ออกแบบ Prompt
          ออกแบบและดำเนินการทดลองเพื่อประเมินผล
          ใช้ความรู้ด้าน Machine Learning, Natural Language Processing และการออกแบบการทดลอง
    \item \textbf{นายบูรณิน บุณโยประการ} (Data Processing) --
          จัดการชุดข้อมูล การ Replay Logs เข้า Splunk การ Normalize
          และการพัฒนาวิธี Event Selection และ Representation
          ใช้ความรู้ด้านระบบปฏิบัติการ ความมั่นคงปลอดภัยไซเบอร์ โครงสร้างข้อมูลและอัลกอริทึมบนกราฟ
\end{enumerate}

\section{\ifenglish%
Impacts of this project on society, health, safety, legal, and cultural issues
\else%
ผลกระทบด้านสังคม สุขภาพ ความปลอดภัย กฎหมาย และวัฒนธรรม
\fi}

\subsection{ผลกระทบด้านสังคม}
โครงงานนี้ช่วยให้หน่วยงานที่มีบุคลากรด้านความปลอดภัยจำกัด สามารถสืบสวนเหตุการณ์ได้รวดเร็วขึ้น
และนำ LLM มาใช้ได้ด้วยต้นทุนที่ต่ำลง ส่งผลให้ระบบสารสนเทศขององค์กรและผู้ใช้บริการมีความปลอดภัยมากขึ้น

\subsection{ผลกระทบด้านสุขภาพ}
การลดปริมาณงานคัดกรอง Logs ที่ซ้ำซากช่วยลดภาระงานและความเหนื่อยล้าจากการแจ้งเตือนจำนวนมาก (Alert Fatigue)
ของ SOC Analyst ทางอ้อม

\subsection{ผลกระทบด้านความปลอดภัย}
ระบบใช้ LLM ที่ทำงานภายในเครื่อง ทำให้ Logs ซึ่งอาจมีข้อมูลสำคัญขององค์กรไม่ต้องถูกส่งออกไปยังผู้ให้บริการภายนอก
อย่างไรก็ดี ผลการวิเคราะห์ของ LLM อาจผิดพลาดได้ ระบบจึงออกแบบให้แสดงหลักฐาน (Evidence) ประกอบทุกคำตอบ
และ SOC Analyst ยังเป็นผู้ตัดสินใจขั้นสุดท้าย

\subsection{ผลกระทบด้านกฎหมาย}
โครงงานใช้เฉพาะชุดข้อมูลสาธารณะที่เผยแพร่เพื่อการวิจัย
และการประมวลผลภายในเครื่องช่วยสนับสนุนการปฏิบัติตามกฎหมายคุ้มครองข้อมูลส่วนบุคคล (PDPA)
เมื่อนำไปใช้กับข้อมูลจริงขององค์กร

\subsection{ผลกระทบด้านวัฒนธรรม}
ไม่มีผลกระทบโดยตรงต่อวัฒนธรรม แต่ส่งเสริมวัฒนธรรมการทำงานด้านความปลอดภัยที่ตัดสินใจโดยอาศัยหลักฐาน (Evidence-based)
